% Part: turing-machines
% Chapter: machines-computations
% Section: church-turing-thesis

\documentclass[../../../include/open-logic-section]{subfiles}

\begin{document}

\olfileid[fr]{tur}{mac}{ctt}
\olsection{La thèse de Church--Turing}

Les machines de Turing visent à donner un sens précis à la
notion intuitive de procédure effective. Turing pensait que
quiconque comprend à la fois cette notion et le fonctionnement
d'une machine de Turing reconnaîtrait intuitivement que tout
ce qu'une procédure effective peut accomplir peut aussi
l'être par une telle machine. Les autres modèles précis
proposés pour formaliser les procédures effectives appuient
cette idée : ils sont extensionnellement équivalents aux
machines de Turing, c'est-à-dire qu'ils calculent exactement
les mêmes fonctions. Cette affirmation est la
\emph{thèse de Church--Turing}. Elle relie une notion
intuitive à un modèle mathématique précis; ce n'est donc
pas un théorème obtenu à partir de deux définitions formelles.

\begin{defn}[Thèse de Church--Turing]
La \emph{thèse de Church--Turing} affirme que tout ce
qui est calculable par une procédure effective est
calculable par une machine de Turing.
\end{defn}

On fait appel à la thèse de Church--Turing de deux façons.
La première sert parfois de prétexte à la paresse.
Supposons qu'une procédure effective permettant de calculer
une fonction soit décrite, par exemple en « pseudocode ».
La thèse permet alors d'affirmer qu'une machine de Turing
calcule la même fonction, même si nous ne l'avons pas
effectivement construite.

Le second usage est plus intéressant du point de vue
philosophique. On peut démontrer que certaines fonctions
ne sont calculables par aucune machine de Turing. La thèse
de Church--Turing permet alors de conclure qu'aucune
procédure effective, quelle qu'elle soit, ne peut les
calculer. En effet, s'il existait une telle procédure,
la thèse entraînerait l'existence d'une machine de Turing
qui calcule la même fonction. Ainsi, démontrer qu'aucune
machine de Turing ne calcule une fonction exclut aussi
toute procédure effective pour la calculer. On invoque
notamment la thèse pour soutenir que le problème dit
de l'arrêt n'est pas seulement insoluble par machine
de Turing : il n'admet aucune solution effective.

\end{document}
